\documentclass[9.5pt,a4paper]{article}
\usepackage[utf8]{inputenc}
\usepackage[T1]{fontenc}
\usepackage{lmodern}
\usepackage[margin=2.0cm]{geometry}
\usepackage{amsmath}
\usepackage{graphicx}
\usepackage{booktabs}
\usepackage{longtable}
\usepackage{array}
\usepackage{caption}
\usepackage{float}
\usepackage{hyperref}
\usepackage{parskip}
\captionsetup{font=small,labelfont=bf}
\newcolumntype{P}[1]{>{\raggedright\arraybackslash}p{#1}}
\newcommand{\file}[1]{\texttt{\footnotesize #1}}
\newcommand{\yes}{\textbf{supported}}
\newcommand{\partly}{\textbf{partly}}
\newcommand{\no}{\textbf{not tested}}
\input{values}

\title{What was done, and the evidence for it\\[2pt]
\large charge states from U-Net line maps}
\author{\small Generated by \texttt{docs/charge\_states/make\_evidence.py}\\
\small run: \texttt{\vRun}}
\date{}

\begin{document}
\maketitle

\noindent
This document lists every step that was done, then every claim made in the
other PDFs, with the number that supports it, the file it is read from, and
how strong the evidence is.  All paths are relative to
\file{results/\vRun/} unless they start with \file{src/}, \file{scripts/},
\file{tests/} or \file{docs/}.  All results are on simulated, noise-free
devices: \vNT{} held-out test devices, and \vNV{} validation devices used
only to choose settings.

\section{What was done, in order}

\begin{longtable}{P{0.4cm}P{6.6cm}P{4.6cm}P{4.6cm}}
\toprule
\# & step & code & output \\
\midrule
\endhead
1 & Checked the sign convention on the true charge states: holes, so
$(0,0)$ at the top right; occupation never grows towards positive voltage.
& \file{study/charge\_states.py} (\texttt{device\_variation}) &
\file{charge\_states/variation\_*.md} \\
2 & Wrote \texttt{label\_states}: regions $\to$ border slopes $\to$ which
side gains $\to$ count from $(0,0)$ along a maximum spanning tree;
consistency check on the remaining borders. &
\file{src/csdrecon/study/charge\_states.py} & -- \\
3 & Tested it on the true line maps; found and fixed a bug (lines stop
1\,px short of the frame). & \texttt{close\_frame}; \file{tests/} & -- \\
4 & Scored every U-Net (15 budgets) on the whole window. &
\file{scripts/run\_16\_charge\_states.py} &
\file{charge\_states/summary.*}, \file{per\_device.csv} \\
5 & Added the conventional route on the full image: three detectors,
\vNDenseSet{} settings scanned on validation; the step-residual detector
was chosen. & \texttt{dense\_line\_map}; run\_16 &
\file{charge\_states/dense\_scan.csv} \\
6 & Restricted the evaluation to the few-electron region ($n\le3$,
absolute labels); defined balanced accuracy and $(1,1)$ landing. &
\texttt{score\_few} & -- \\
7 & Added the line fit (straight segments, free ends extended to the next
line or the frame); \vNFitSet{} settings scanned on validation. &
\texttt{fit\_lines}; \file{scripts/run\_17\_few\_electron.py} &
\file{few\_electron/fit\_scan.csv}, \file{summary.*} \\
8 & Recomputed the line-level metrics from the checkpoints actually used
(some budgets were retrained with another seed). &
\file{docs/charge\_states/make\_eval.py} &
\file{evaluation\_note/} \file{line\_metrics.csv} \\
9 & Bootstrap 95\,\% intervals over the test devices. & make\_eval,
make\_report & \file{evaluation\_note/ci\_table.tex} \\
10 & Per-device F1@1 against charge-state accuracy. & make\_report &
\file{report\_note/figures/} \file{f1\_vs\_states.pdf} \\
11 & Equal-cost comparison: rays + interpolation, coarse scan upsampled,
coarse scan native; each tuned on validation; paired bootstrap against
U-Net + fit. & \file{scripts/run\_18\_equal\_cost.py} &
\file{equal\_cost/scan\_*.csv}, \file{summary.*}, \file{per\_device.csv} \\
12 & Wrote the PDFs listed in Sec.~\ref{sec:files}. & \file{docs/charge\_states/} & see Sec.~\ref{sec:files} \\
\bottomrule
\end{longtable}

\section{Claims and their evidence}

Each claim is marked \yes{} (direct measurement, with an interval where it
matters), \partly{} (true for what was tested, but a stronger test is still
missing) or \no.  Numbers are means over the \vNT{} test devices; the
few-electron region is $n_1,n_2\le\vNmax$.

\begin{longtable}{P{5.2cm}P{4.7cm}P{4.4cm}P{1.7cm}}
\toprule
claim & evidence & source file & status \\
\midrule
\endhead

\multicolumn{4}{l}{\textit{The labelling method}} \\[2pt]
The sign convention is as assumed: corner $(0,0)$, holes. &
corner $(0,0)$ on \vCornerT/\vNT{} test and \vCornerV/\vNV{} validation
devices; $n_1$, $n_2$ never grow towards $+V$ on all of them. &
\file{charge\_states/per\_device.csv} (\texttt{corner\_is\_00},
\texttt{n*\_grows\_to\_neg\_V*}) & \yes \\
\texttt{label\_states} is correct on complete line maps. &
true maps: pixel accuracy \vFullAccRel\,\% (test), \vFullAccRelVal\,\%
(validation); state recall \vFullRecall\,\%; few-electron balanced
\vFullBalFew\,\%, landing \vFullLandFew\,\%; consistency \vFullCons. &
\file{charge\_states/summary*.csv}, \file{few\_electron/summary.csv} &
\yes \\
The method is tested in code. & \vTestsCharge{} unit tests for
labelling, line fit and scoring; whole suite \vTestsPassed{} passed,
\vTestsFailed{} failed (run when this PDF was built). &
\file{tests/test\_charge\_states.py} & \yes \\
\midrule

\multicolumn{4}{l}{\textit{The U-Net maps}} \\[2pt]
The U-Net finds the lines (line level). & F1@1 =
\vLineFoneMin--\vLineFoneMax{} over the 15 budgets (\vLineFone{} at
$8\times60$). & \file{evaluation\_note/} \file{line\_metrics.csv} & \yes \\
But from its maps alone few charge states are right. & balanced few-e
\vUfourBal\,\% ($4\times40$) to \vUBal\,\% [\vUBalLo, \vUBalHi]
($8\times60$); landing \vUfourLand--\vULand\,\%. &
\file{few\_electron/summary.csv}, \file{per\_device.csv} & \yes \\
Line-level F1 does not predict charge-state usefulness. & Spearman
$\rho = \vRhoU$ (per device, $8\times60$); devices with F1@1 $\ge0.85$
(\vNHiFone) range \vHiFoneBalMin--\vHiFoneBalMax\,\% balanced. &
\file{report\_note/figures/} \file{f1\_vs\_states.pdf} & \yes{} (one budget) \\
The failure is gaps in the lines. & seen on every inspected failure;
the line fit, which only closes gaps, raises balanced accuracy at every
budget. & figures in \file{charge\_states\_method.pdf},
\file{few\_electron\_results.pdf} & \partly{} (no automatic gap count) \\
\midrule

\multicolumn{4}{l}{\textit{The line fit}} \\[2pt]
The line fit improves the U-Net maps. & +\vGainMin{} to +\vGainMax{}
points balanced at all \vNBud{} budgets; $8\times60$: \vUBal{}
$\to$ \vFBal\,\% [\vFBalLo, \vFBalHi]; validation \vFitValWithout{} $\to$
\vFitValWith\,\%. & \file{few\_electron/summary.csv},
\file{fit\_scan.csv} & \yes \\
It can also draw false lines. & visible on sample\_223 (extension through
the $(1,0)$ region). & \file{few\_electron/figures/} \file{fit\_steps.png} &
\partly{} (not counted) \\
\midrule

\multicolumn{4}{l}{\textit{Against the conventional route}} \\[2pt]
On the full image, the conventional route is far better. & balanced
\vDBal\,\% [\vDBalLo, \vDBalHi], landing \vDLand\,\%, with
\vFold$\times$ more measurements. & \file{few\_electron/summary.csv} &
\yes \\
At equal measurement cost, U-Net + fit beats the conventional route. &
\vEqWins/\vEqNBud{} budgets with paired 95\,\% interval above 0;
$8\times60$: \vFBal{} vs \vEqCompBal\,\%, difference \vEqDiff{}
[\vEqDiffLo, \vEqDiffHi]. & \file{equal\_cost/per\_device.csv},
\file{summary.csv} & \partly{} (simple competitors only) \\
The competitors were not handicapped by tuning. & validation best: rays
\vScanRaysBest\,\% (\vScanRaysN{} settings), coarse
\vScanCoarseBest\,\% (\vScanCoarseN), native \vScanNativeBest\,\%
(\vScanNativeN); a wider threshold range changed little. &
\file{equal\_cost/scan\_*.csv} & \partly{} (see Sec.~\ref{sec:checks}) \\
Without the fit the U-Net is not always better. & raw U-Net above the
best competitor's mean at \vEqRawWins/\vEqNBud{} budgets. &
\file{equal\_cost/summary.csv} & \yes \\
\midrule

\multicolumn{4}{l}{\textit{Other}} \\[2pt]
Fewer honeycombs do not rescue the (current) U-Net. & state recall
similar in every density group; pixel accuracy rises with the
baseline. & \file{charge\_states/results\_note/} \file{table\_A\_density.tex} &
\partly{} (no retrained U-Net) \\
The U-Net is useful for tuning to $(1,1)$. & landing at most
\vBestLand\,\% even with the fit. & \file{few\_electron/summary.csv} &
\textbf{not supported} \\
It works with noise / on real data / against the earlier reconstruction. &
-- & -- & \no \\
\bottomrule
\end{longtable}

\section{Checks that were run, and mistakes found}
\label{sec:checks}

\begin{longtable}{P{5.0cm}P{11.2cm}}
\toprule
check or mistake & what happened \\
\midrule
\endhead
Labelling on the true maps & First version: below 40\,\% accuracy on the
first training devices checked (development observation, not stored).
Cause: lines stop 1\,px short of the frame. Fixed with
\texttt{close\_frame}; now \vFullAccRel\,\%. \\
Dense detector on the full image & Canny and gradient thresholds drew the
sensor background at the frame, not the transitions; replaced by the
step-residual detector, chosen on validation. \\
Fig.~2 error-map legend & The first legend described the colours wrongly;
corrected and re-rendered. \\
Negative labels in figures & Labels such as $(8,-1)$ were drawn as
``7,999''; display fixed (metrics were not affected). \\
Checkpoint mismatch & Some \file{evaluation/results.txt} files describe
seed-0 models; the charge-state studies use the retrained seed. Line metrics
were recomputed from the used checkpoints (e.g.\ $8\times60$:
\file{model\_seed2}). \\
Competitor handicap (equal cost) & First run: threshold range too narrow,
coarse scan only upsampled. Re-run with a wider range and a native coarse
variant; the result held. \\
Unfair-advantage sanity check & The coarse route tuned on two test devices
themselves reached only 24\,\% and 14\,\% (a check run during the analysis;
not stored in a file). \\
``Best competitor'' choice & Picked per budget with hindsight on the test
set, which favours the competitor, not the U-Net. \\
A folder created by mistake & A run created an empty
\file{results/4-6-8-12-16-24\_.../charge\_states}; it was removed and the
script fixed. \\
Unit tests & \vTestsPassed{} passed, \vTestsFailed{} failed. \\
\bottomrule
\end{longtable}

\section{Key evidence, at a glance}

\begin{table}[H]
\centering
\caption{Equal measurement cost: U-Net + fit against the best conventional
variant at each budget (balanced few-electron accuracy, \%).  Source:
\file{equal\_cost/per\_device.csv}.}
\small
\input{table_equal.tex}
\end{table}

\begin{figure}[H]
\centering
\includegraphics[width=0.95\textwidth]{figures/equal_cost_curves.pdf}
\caption{Every method, every budget.  Source:
\file{equal\_cost/summary.csv}.}
\end{figure}

\begin{figure}[H]
\centering
\includegraphics[width=0.5\textwidth]{figures/f1_vs_states.pdf}
\caption{Line-level F1@1 against charge-state accuracy, per test device,
$8\times60$.  Source: recomputed in \file{docs/charge\_states/make\_report.py}
from \file{few\_electron/per\_device.csv} and the U-Net checkpoint.}
\end{figure}

\section{What is not evidenced}
\begin{itemize}
\item Any result with sensor noise, or on an experimental CSD.
\item A comparison with the project's earlier probabilistic reconstruction
  (or any published reconstruction method) at equal cost.
\item The 4--24-ray sweep: its results folder is missing from
  \file{results/}.
\item U-Nets retrained on a narrower, experimentally realistic capacitance
  range, or with a loss that punishes gaps.
\item Novelty: no literature check has been made.
\end{itemize}

\section{Files}
\label{sec:files}
\begin{longtable}{P{5.6cm}P{10.6cm}}
\toprule
file & content \\
\midrule
\endhead
\file{charge\_states\_report.pdf} & the full account: question, method,
results, interpretation, contribution \\
\file{charge\_states\_method.pdf} & the labelling method, step by step \\
\file{charge\_states\_results.pdf} & whole-window results, dense
reference, honeycomb density \\
\file{few\_electron\_results.pdf} & few-electron region and the line fit \\
\file{evaluation\_method.pdf} & how accuracy is evaluated \\
\file{evidence.pdf} & this document \\
\file{charge\_states/} & run\_16 tables, per-device data, dense scan \\
\file{few\_electron/} & run\_17 tables, per-device data, fit scan \\
\file{equal\_cost/} & run\_18 tables, per-device data, competitor scans \\
\file{src/csdrecon/study/charge\_states.py} & the library: labelling,
line fit, detectors, metrics, figures \\
\file{scripts/run\_16..18\_*.py} & the three studies \\
\file{docs/charge\_states/make\_*.py, *.tex} & the PDF builders \\
\bottomrule
\end{longtable}

\end{document}
